% Préambule « think tank » de la note v3 (sans empattement, bandeau, encadrés, chiffres clés).
\usepackage[T1]{fontenc}
\usepackage[utf8]{inputenc}
\usepackage[scaled=0.95]{helvet}
\renewcommand{\familydefault}{\sfdefault}
\usepackage[french]{babel}
\usepackage[a4paper,left=2.0cm,right=2.0cm,top=1.7cm,bottom=1.8cm]{geometry}
\usepackage{graphicx}
\usepackage{booktabs,tabularx,array}
\usepackage[table]{xcolor}
\usepackage{titlesec}
\usepackage[labelfont={bf,color=bleu},labelsep=period,font={small,sf},singlelinecheck=false,justification=raggedright,skip=3pt]{caption}
\usepackage{fancyhdr}
\usepackage{enumitem}
\usepackage{microtype}
\usepackage{textcomp}
\usepackage{amsmath}
\usepackage[most]{tcolorbox}
\providecommand{\euro}{\texteuro}
\usepackage[hidelinks]{hyperref}

\definecolor{encre}{HTML}{1A1A1A}
\definecolor{bleu}{HTML}{173F8A}
\definecolor{bleuclair}{HTML}{E8EEF8}
\definecolor{brique}{HTML}{B8642A}
\definecolor{gris}{HTML}{5A5A5A}

\titleformat{\section}{\normalfont\Large\bfseries\color{bleu}}{}{0em}{}[{\color{brique}\titlerule[1.2pt]}]
\titleformat{\subsection}{\normalfont\normalsize\bfseries\color{bleu}}{}{0em}{}
\titlespacing*{\section}{0pt}{2.2ex plus .6ex minus .2ex}{0.9ex plus .2ex}
\titlespacing*{\subsection}{0pt}{1.2ex plus .5ex minus .2ex}{0.3ex plus .1ex}

\pagestyle{fancy}
\fancyhf{}
\renewcommand{\headrulewidth}{0pt}
\fancyfoot[R]{\small\color{gris}\thepage}
\fancyfoot[L]{\small\color{gris}Taxer mieux plutôt que davantage}

\setlength{\parindent}{0pt}
\setlength{\parskip}{0.5em}
\linespread{1.04}
\setlist{itemsep=1pt,topsep=1pt,parsep=0pt,leftmargin=1.2em}
\setlength{\textfloatsep}{8pt plus 2pt minus 2pt}
\setlength{\floatsep}{6pt plus 2pt minus 2pt}
\renewcommand{\topfraction}{0.88}\renewcommand{\bottomfraction}{0.7}
\renewcommand{\textfraction}{0.1}\renewcommand{\floatpagefraction}{0.8}

% Encadré « en bref » et chiffres clés
\newtcolorbox{enbref}[1][]{colback=bleuclair,colframe=bleuclair,boxrule=0pt,arc=0pt,left=8pt,right=8pt,top=6pt,
  bottom=6pt,fontupper=\small,#1}
\newtcolorbox{exemple}[1]{colback=white,colframe=brique,boxrule=0.6pt,leftrule=3pt,arc=0pt,left=6pt,right=6pt,top=4pt,
  bottom=4pt,fontupper=\small,title={\bfseries #1},coltitle=brique,colbacktitle=white,attach title to upper={\ \ }}
\newcommand{\chiffre}[2]{\begin{minipage}[t]{0.235\linewidth}\centering{\Large\bfseries\color{bleu}#1}\\[1pt]
  {\scriptsize\color{gris}#2}\end{minipage}}

\newcommand{\graphique}[3]{%
  \begin{figure}[!htbp]\centering
  \caption{#2}\label{#3}\vspace{1pt}
  \includegraphics[width=\linewidth]{figures/#1.pdf}\par\vspace{1pt}
  {\raggedright\scriptsize\color{gris}\input{figures/#1_source.tex}\par}
  \end{figure}}
\newcommand{\deuxgraphiques}[6]{%
  \begin{figure}[!htbp]
  \begin{minipage}[t]{0.485\linewidth}\caption{#2}\label{#3}\vspace{1pt}
  \includegraphics[width=\linewidth]{figures/#1.pdf}\par
  {\raggedright\scriptsize\color{gris}\input{figures/#1_source.tex}\par}\end{minipage}\hfill
  \begin{minipage}[t]{0.485\linewidth}\caption{#5}\label{#6}\vspace{1pt}
  \includegraphics[width=\linewidth]{figures/#4.pdf}\par
  {\raggedright\scriptsize\color{gris}\input{figures/#4_source.tex}\par}\end{minipage}
  \end{figure}}

\DeclareUnicodeCharacter{2212}{\ensuremath{-}}
\DeclareUnicodeCharacter{2248}{\ensuremath{\approx}}
\DeclareUnicodeCharacter{00D7}{\ensuremath{\times}}
\DeclareUnicodeCharacter{2192}{\ensuremath{\rightarrow}}
\DeclareUnicodeCharacter{2264}{\ensuremath{\leq}}
\DeclareUnicodeCharacter{2265}{\ensuremath{\geq}}
\DeclareUnicodeCharacter{202F}{\,}
\DeclareUnicodeCharacter{2009}{\,}
\DeclareUnicodeCharacter{03C3}{\ensuremath{\sigma}}
\DeclareUnicodeCharacter{03C4}{\ensuremath{\tau}}
\DeclareUnicodeCharacter{03C1}{\ensuremath{\rho}}
\DeclareUnicodeCharacter{03B5}{\ensuremath{\varepsilon}}
\DeclareUnicodeCharacter{03B8}{\ensuremath{\theta}}
\DeclareUnicodeCharacter{03B4}{\ensuremath{\delta}}
\DeclareUnicodeCharacter{03BB}{\ensuremath{\lambda}}
\DeclareUnicodeCharacter{00B2}{\textsuperscript{2}}
